\pdfbookmark{Общая характеристика работы}{characteristic}             % Закладка pdf
\section*{Общая характеристика работы}

\newcommand{\actuality}{\pdfbookmark[1]{Актуальность}{actuality}\underline{\textbf{\actualityTXT}}}
\newcommand{\progress}{\pdfbookmark[1]{Разработанность темы}{progress}\underline{\textbf{\progressTXT}}}
\newcommand{\aim}{\pdfbookmark[1]{Цели}{aim}\underline{{\textbf\aimTXT}}}
\newcommand{\tasks}{\pdfbookmark[1]{Задачи}{tasks}\underline{\textbf{\tasksTXT}}}
\newcommand{\aimtasks}{\pdfbookmark[1]{Цели и задачи}{aimtasks}\aimtasksTXT}
\newcommand{\novelty}{\pdfbookmark[1]{Научная новизна}{novelty}\underline{\textbf{\noveltyTXT}}}
\newcommand{\influence}{\pdfbookmark[1]{Практическая значимость}{influence}\underline{\textbf{\influenceTXT}}}
\newcommand{\methods}{\pdfbookmark[1]{Методология и методы исследования}{methods}\underline{\textbf{\methodsTXT}}}
\newcommand{\defpositions}{\pdfbookmark[1]{Положения, выносимые на защиту}{defpositions}\underline{\textbf{\defpositionsTXT}}}
\newcommand{\reliability}{\pdfbookmark[1]{Достоверность}{reliability}\underline{\textbf{\reliabilityTXT}}}
\newcommand{\probation}{\pdfbookmark[1]{Апробация}{probation}\underline{\textbf{\probationTXT}}}
\newcommand{\contribution}{\pdfbookmark[1]{Личный вклад}{contribution}\underline{\textbf{\contributionTXT}}}
\newcommand{\publications}{\pdfbookmark[1]{Публикации}{publications}\underline{\textbf{\publicationsTXT}}}

{\actuality} Одним из важных компонентов динамического тренажера автотранспорта и спецтехники (АТ и СТ) является система подвижности (СП), предназначенная для выработки у обучаемого необходимых моторных навыков в процессе управления техникой. Применяемые в большинстве тренажеров системы подвижности представляют собой манипулятор с параллельной кинематикой, состоящий из подвижной платформы, связанной с неподвижным основанием одним или несколькими поступательными звеньями. В зависимости от требований к имитируемым эффектам, используются платформы с числом степеней свободы от 1 до 6. В подавляющем большинстве тренажеров АТ и СТ (порядка 70-80~\% продукции ведущих производителей тренажерных систем) применяются системы подвижности 2DOF и 3DOF с кривошипно-шатунным механизмом (КШМ) привода звеньев, обеспечивающие отработку необходимых при имитации данной техники координат подвижной платформы (крен, тангаж, высота). При разработке ПО управления универсальной системой подвижности (УСП) необходимо учитывать такие факторы, как геометрические параметры платформы, рабочую зону манипулятора, массогабаритные характеристики платформы в различных состояниях, алгоритмы управления положением привода, протоколы обмена с датчиками и исполнительными устройствами. При наличии в номенклатуре производимых тренажеров нескольких систем подвижности с различными кинематическими схемами и типами приводов достаточно остро встает вопрос о разработке ПО, позволяющего осуществлять управление несколькими СП по стандартизованному протоколу обмена и производить быструю настройку под конкретную СП с возможностью диагностики состояния оборудования.	
 
 \textbf{Актуальность} рассматриваемой диссертации определяется потребностями:
\begin{itemize}
	\item создания систем подвижности, обеспечивающих наиболее реалистичные эффекты движения имитируемой техники;
	\item обеспечения безопасной эксплуатации АРМ обучаемого как подвижного объекта с людьми, что требует применения специальных мер (дополнительные датчики, дублирование систем);
	\item унификации модулей тренажерного комплекса, в частности, системы подвижности, что позволяет существенно упростить процесс развертывания тренажера в ходе пуско-наладочных работ.
В работе рассматриваются вопросы реализации системы управления несколькими типами СП, построение архитектуры приложения, отвечающей заданным требованиям универсализации, а также применяемые при этом инженерные решения в области аппаратной части платформы.
\end{itemize}
% {\progress} 
% Этот раздел должен быть отдельным структурным элементом по
% ГОСТ, но он, как правило, включается в описание актуальности
% темы. Нужен он отдельным структурынм элемементом или нет ---
% смотрите другие диссертации вашего совета, скорее всего не нужен.

{\aim} данной работы является разработка унифицированной системы управления динамической платформой подвижности, позволяющей абстрагировать реализацию алгоритмов управления и ограничений, накладываемых рабочей зоной платформы, от клиентской стороны -- модуля тренажерного комплекса, отвечающего за преобразование и передачу данных телеметрии от математической модели техники к системе подвижности

Для~достижения поставленной цели необходимо было решить следующие {\tasks}:
\begin{enumerate}
  \item разработать алгоритмы управления СП, учитывающие динамические характеристики подвижной платформы и закрепленного на ней объекта (соответствует п. 5 паспорта научной специальности);
  \item разработать методику настройки регуляторов СП, обеспечивающие требуемое качество процессов в системе управления (соответствует п. 4 паспорта научной специальности);
  \item разработать алгоритмы формирования траектории в соответствии с рассчитанной рабочей зоной платформы (соответствует п. 6 паспорта научной специальности);
  \item произвести параметрическую идентификацию модели СП по данным от датчиков (соответствует п. 8 паспорта научной специальности);
  \item разработать процедуры автоматической диагностики состояния оборудования СП (соответствует п. 15 паспорта научной специальности);
  \item разработать алгоритмы взаимодействия клиенткой стороны и контроллера движения СП (соответствует п. 17 паспорта научной специальности).
\end{enumerate}

{\novelty}
\begin{enumerate}
  \item получено аналитическое решение уравнений динамики трехстепенной электромеханической платформы с кривошипно-шатунных звеньями и пружинной опорой, пригодное для расчета в режиме реального времени, что позволяет осуществить коррекцию траектории. Это значительно улучшает отклик системы подвижности и уменьшает ошибку регулирования;
  \item разработан алгоритм коррекции траектории движения УСП, учитывающий кинематику и динамику механизма в совокупности с ограничениями приводных систем. Алгоритм предполагает "мягкую" срезку заданных перемещений, скоростей и ускорений, что обеспечивает наилучшее приближение реальной траектории платформы к заданной;
  \item получено аналитическое решение прямой задачи кинематики, позволяющее с достаточной точностью вычислять в режиме реального времени текущие угловые и поступательные координаты СП без применения итерационных алгоритмов, что позволяет не только визуализировать текущее положение платформы без использования дополнительных датчиков, но и применять продвинутые методики управления с обратной связью по координатам платформы;
  \item проведено исследование точности позиционирования СП при помощи прецизионного промышленного датчика углов наклона, экспериментально подтвердившее обоснованность использования описанного выше решения прямой задачи кинематики в задачах управления платформой.
\end{enumerate}

{\influence} определяется внедрением ПАК управления системой подвижности в тренажерных комплексах военной, гражданской техники и речных судов производства АО «Кронштадт Технологии». Результаты работы могут быть использованы при разработке систем управления подвижными платформами тренажеров автотранспорта и спецтехники.

{\methods} Для решения задач, рассматриваемых в работе использовались методы теории систем автоматического управления, линейной алгебры, вычислительной математики. Для создания программно-аппаратного комплекса применены методы технологии программирования. В ходе верификации результатов моделирования использовались методы планирования вычислительных экспериментов и методы обработки данных натурных испытаний.

{\defpositions}
\begin{itemize}
  \item упрощенная модель динамики трехстепенной динамической платформы на кривошипно-шатунных опорах, позволяющая с приемлемой точностью рассчитывать вращающие моменты двигателей по заданным перемещениям, скоростям и ускорениям СП в режиме реального времени; 
  \item алгоритм управления координатами УСП, учитывающий динамику механизма и приводных систем;
  \item алгоритм ограничения координат и безударного выхода на границу рабочей зоны УСП;
  \item методика настройки параметров регулятора с целью достижения заданных требований к УСП;
\end{itemize}

{\reliability} Достоверность исследований, проведенных в диссертации, подтверждается:
\begin{itemize}
	\item проверкой теоретических выкладок и уравнений с помощью наборов входных тестовых данных;
	\item проведенными экспериментальными исследованиями с паспортизацией предельных достижимых характеристик УСП;
	\item тестированием работы УСП специалистами Компании АО «Кронштадт Технологии».
\end{itemize}

{\probation}

Основные результаты диссертационной работы докладывались и обсуждались на конференциях:
\begin{itemize}
	\item на конференции VII Поляховские Чтения – г. Санкт-Петербург, февраль 2015 г.
	\item на конференции Сетевое партнерство в науке, промышленности и образовании – г. Санкт-Петербург, июль 2016 г.,
\end{itemize}

выставках:
\begin{itemize}
	\item Строительная Техника и Технологии (СТТ), 2015, г. Москва,
	\item Международный военно-технический форум «АРМИЯ-2015», 2015 г., г. Кубинка, Московская область,
	\item Ежегодное совещание командования инженерных войск по округам, 2015 г., г. Кстово, Нижегородская область ;
\end{itemize}

\ifnumequal{\value{bibliosel}}{0}
{%%% Встроенная реализация с загрузкой файла через движок bibtex8. (При желании, внутри можно использовать обычные ссылки, наподобие `\cite{vakbib1,vakbib2}`).
    {\publications} Основные результаты по теме диссертации изложены
    в~XX~печатных изданиях,
    X из которых изданы в журналах, рекомендованных ВАК,
    X "--- в тезисах докладов.
}%
{%%% Реализация пакетом biblatex через движок biber
    \begin{refsection}[bl-author, bl-registered]
        % Это refsection=1.
        % Процитированные здесь работы:
        %  * подсчитываются, для автоматического составления фразы "Основные результаты ..."
        %  * попадают в авторскую библиографию, при usefootcite==0 и стиле `\insertbiblioauthor` или `\insertbiblioauthorgrouped`
        %  * нумеруются там в зависимости от порядка команд `\printbibliography` в этом разделе.
        %  * при использовании `\insertbiblioauthorgrouped`, порядок команд `\printbibliography` в нём должен быть тем же (см. biblio/biblatex.tex)
        %
        % Невидимый библиографический список для подсчёта количества публикаций:
        \phantom{\printbibliography[heading=nobibheading, section=1, env=countauthorvak,          keyword=biblioauthorvak]%
        \printbibliography[heading=nobibheading, section=1, env=countauthorwos,          keyword=biblioauthorwos]%
        \printbibliography[heading=nobibheading, section=1, env=countauthorscopus,       keyword=biblioauthorscopus]%
        \printbibliography[heading=nobibheading, section=1, env=countauthorconf,         keyword=biblioauthorconf]%
        \printbibliography[heading=nobibheading, section=1, env=countauthorother,        keyword=biblioauthorother]%
        \printbibliography[heading=nobibheading, section=1, env=countregistered,         keyword=biblioregistered]%
        \printbibliography[heading=nobibheading, section=1, env=countauthorpatent,       keyword=biblioauthorpatent]%
        \printbibliography[heading=nobibheading, section=1, env=countauthorprogram,      keyword=biblioauthorprogram]%
        \printbibliography[heading=nobibheading, section=1, env=countauthor,             keyword=biblioauthor]%
        \printbibliography[heading=nobibheading, section=1, env=countauthorvakscopuswos, filter=vakscopuswos]%
        \printbibliography[heading=nobibheading, section=1, env=countauthorscopuswos,    filter=scopuswos]}%
        %
        \nocite{*}%
        %
        {\publications} Основные результаты по теме диссертации изложены в~\arabic{citeauthor}~печатных изданиях,
        \arabic{citeauthorvak} из которых изданы в журналах, рекомендованных ВАК%
        \ifnum \value{citeauthorscopuswos}>0%
            , \arabic{citeauthorscopuswos} "--- в~периодических научных журналах, индексируемых Web of~Science и Scopus%
        \fi%
        \ifnum \value{citeauthorconf}>0%
            , \arabic{citeauthorconf} "--- в~тезисах докладов.
        \else%
            .
        \fi%
        \ifnum \value{citeregistered}=1%
            \ifnum \value{citeauthorpatent}=1%
                Зарегистрирован \arabic{citeauthorpatent} патент.
            \fi%
            \ifnum \value{citeauthorprogram}=1%
                Зарегистрирована \arabic{citeauthorprogram} программа для ЭВМ.
            \fi%
        \fi%
        \ifnum \value{citeregistered}>1%
            Зарегистрированы\ %
            \ifnum \value{citeauthorpatent}>0%
            \formbytotal{citeauthorpatent}{патент}{}{а}{}%
            \ifnum \value{citeauthorprogram}=0 . \else \ и~\fi%
            \fi%
            \ifnum \value{citeauthorprogram}>0%
            \formbytotal{citeauthorprogram}{программ}{а}{ы}{} для ЭВМ.
            \fi%
        \fi%
        % К публикациям, в которых излагаются основные научные результаты диссертации на соискание учёной
        % степени, в рецензируемых изданиях приравниваются патенты на изобретения, патенты (свидетельства) на
        % полезную модель, патенты на промышленный образец, патенты на селекционные достижения, свидетельства
        % на программу для электронных вычислительных машин, базу данных, топологию интегральных микросхем,
        % зарегистрированные в установленном порядке.(в ред. Постановления Правительства РФ от 21.04.2016 N 335)
    \end{refsection}%
    \begin{refsection}[bl-author, bl-registered]
        % Это refsection=2.
        % Процитированные здесь работы:
        %  * попадают в авторскую библиографию, при usefootcite==0 и стиле `\insertbiblioauthorimportant`.
        %  * ни на что не влияют в противном случае
        \nocite{vakbib2}%vak
        \nocite{patbib1}%patent
        \nocite{progbib1}%program
        \nocite{bib1}%other
        \nocite{confbib1}%conf
    \end{refsection}%
        %
        % Всё, что вне этих двух refsection, это refsection=0,
        %  * для диссертации - это нормальные ссылки, попадающие в обычную библиографию
        %  * для автореферата:
        %     * при usefootcite==0, ссылка корректно сработает только для источника из `external.bib`. Для своих работ --- напечатает "[0]" (и даже Warning не вылезет).
        %     * при usefootcite==1, ссылка сработает нормально. В авторской библиографии будут только процитированные в refsection=0 работы.
}
 % Характеристика работы по структуре во введении и в автореферате не отличается (ГОСТ Р 7.0.11, пункты 5.3.1 и 9.2.1), потому её загружаем из одного и того же внешнего файла, предварительно задав форму выделения некоторым параметрам


%Диссертационная работа была выполнена при поддержке грантов ...

%\underline{\textbf{Объем и структура работы.}} Диссертация состоит из~введения, четырех глав, заключения и~приложения. Полный объем диссертации \textbf{ХХХ}~страниц текста с~\textbf{ХХ}~рисунками и~5~таблицами. Список литературы содержит \textbf{ХХX}~наименование.

%\newpage
\section*{Содержание работы}
Во \underline{\textbf{введении}} обоснована актуальность темы исследования и практическая значимость работы, приведены использованные методы исследования и положения, выносимые на защиту.

В \underline{\textbf{первой главе}} приведен обзор основных средств имитации тактильных воздействий на обучаемого в тренажерах автотранспорта и спецтехники. Рассмотрены типы кинематических схем (механизмов) и приводных систем, применяемых для реализации систем подвижности данного класса тренажеров. Проанализированы характерные акселерационные эффекты движения, для которых приведены графики перемещений и скоростей. Сформулированы требования к приводам системы подвижности с точки зрения возможностей по обеспечению тактильной обратной связи.

% картинку можно добавить так:
%\begin{figure}[ht] 
%  \center
%  \includegraphics [scale=0.27] {latex}
%  \caption{Подпись к картинке.} 
%  \label{img:latex}
%\end{figure}

\underline{\textbf{Вторая глава}} посвящена исследованию кинематических характеристик СП. Рассмотрены уравнения одно-, двух- и трехстепенных платформ с кривошипношатунными опорами. На рисунке \ref{img:3DOFKinematicsScheme} приведена расчетная кинематическая схема трехстепенной платформы.

\begin{figure}[ht!]
	\center
	\includegraphics[scale=0.2]{3dof_kinematics_scheme.png}
	\caption{Расчетная кинематическая схема трехстепенной платформы}
	\label{img:3DOFKinematicsScheme}
\end{figure}

где $\alpha$ -- угол крена, $\beta$ -- угол тангажа, $h$ -- вертикальное смещение геометрического центра подвижной платформы, $q_1,~q_2,~q_3$ -- углы поворота кривошипов, $a,~b,~c,~R,~L$ -- геометрические параметры платформы.

Для данного типа платформ уравнения обратной задачи кинематики, позволяющие рассчитывать углы поворота кривошипов по заданным положениям платформы, выглядят следующим образом:

\begin{equation}
\label{eq:InverseKinematics3DOF}
\begin{aligned}
	q_i & = -2\arctan\frac{\bar{B_i}\pm\sqrt{\bar{B_i^2} + \bar{C_i^2} - \bar{A_i^2}}}{\bar{A_i} - \bar{C_i}} \\
	\bar{A_i} & = (x_{Ci} - x_{Ai})^2 + (y_{Ci} - y_{Ai})^2 + (z_{Ci} - z_{Ai})^2 + R^2 + L^2,\\
	\bar{B_i} & = -2R(y_{Ci} - y_{Ai}),\\
	\bar{C_i} & = 2R(z_{Ci} - z_{Ai}).
\end{aligned}
\end{equation}

где координаты точек $A_i,~B_i,~C_i$ зависят от углов крена, тангажа и подъема платформы.
Для целей последующего анализа динамики СП было рассмотрено решение уравнений прямой задачи кинематики численным и аналитическим способами.

Для решения ПЗК аналитически исходные уравнения для длин кривошипов $\left(x_{Bi}-x_{C_i}\right)^2  + \left(y_{Bi}-y_{C_i}\right)^2 + \left(z_{Bi}-z_{C_i}\right)^2 = L^2, i=1,\ldots,3$ были разложены в ряд Тейлора по степеням показателя $\kappa = R/L$.
 
Были также построены ограничительные кинематические характеристики для двух основных типов трехстепенных платформ: с симметричным и несимметричным расположением центров валов двигателей. Ограничительная характеристика показана на рисунке \ref{img:RestrictKinematicZone}.

\begin{figure}[ht!]
    \center{
        \hfill
        \subcaptionbox[List-of-Figures entry]{\label{img:RestrictKinematicAlphaHNS}} {\includegraphics[width=0.5\linewidth]{RestrictKinematicAlphaHNS}}%
        \hfill       
        \subcaptionbox{\label{img:RestrictKinematicAlphaH}} {\includegraphics[width=0.5\linewidth]{RestrictKinematicAlphaH}}
        \hfill
    }
    \center{
        \subcaptionbox[List-of-Figures entry]{\label{img:RestrictKinematicBetaHNS}} {\includegraphics[width=0.5\linewidth]{RestrictKinematicBetaHNS}}%
        \hfill       
        \subcaptionbox{\label{img:RestrictKinematicBetaH}} {\includegraphics[width=0.5\linewidth]{RestrictKinematicBetaH}}
        \hfill
    }
	\center{    
        \subcaptionbox[List-of-Figures entry]{\label{img:RestrictKinematicAlphaBetaNS}} {\includegraphics[width=0.5\linewidth]{RestrictKinematicAlphaBetaNS}}%
        \hfill       
        \subcaptionbox{\label{img:RestrictKinematicAlphaBeta}} {\includegraphics[width=0.5\linewidth]{RestrictKinematicAlphaBeta}}
        \hfill
    }
    \caption{Проекции зоны кинематических ограничений для симметричной (а,в,д) и несимметричной (б,г,е) СП} % Этот текст попадает в названия рисунков в списке рисунков
    \label{img:RestrictKinematicZone}
\end{figure}

\underline{\textbf{Третья глава}} посвящена исследованию динамики электромеханической системы подвижности на кривошипно-шатунных опорах.

Расчетная схема платформы приведена на рисунке \ref{img:3dof_dynamic_scheme}. Особенностью используемой кинематической схемы является центральный карданный шарнир, удерживающий платформу от перемещений по курсу. 
 
\begin{figure}[ht!]
	\center
	\includegraphics[width=0.65\textwidth]{3dof_kinematics_scheme_dynamic}
	\caption{Расчётная силовая схема подвижности 3DOF}
	\label{img:3dof_dynamic_scheme}
\end{figure}

В ходе работы были получены уравнения динамики, решение которых сводится к нахождению сил реакций шатунов из решения следующей линейной системы уравнений:

\begin{equation}
\label{eq:KoefsLinearSystem2}
\begin{aligned}
& \begin{pmatrix}
F_1 & F_2 & F_3 \\ 
bD_1 & -bD_2 & 0 \\
aD_1 & aD_2 & -cD_3
\end{pmatrix}k + 
\begin{pmatrix}0 \\ S_y^0 \\ -S_x^0\end{pmatrix}kz_c +
\begin{pmatrix}0 \\ -y_c \\ x_c \end{pmatrix}Ek = \frac{1}{L} 
\begin{pmatrix}
m\ddot{h} \\
m_{\alpha} \\
m_{\beta}	
\end{pmatrix} + \\ 
& + \frac{1}{L}\begin{pmatrix}
k_sh - f_0 + mg \\
y_c(m\ddot{h} + mg) \\
-x_c(m\ddot{h} + mg)
\end{pmatrix}
\end{aligned}
\end{equation}

где $k$ -- модули сил реакций шатунов, $x_c,y_c,z_c$ -- положение центра масс в системе координат платформы, $m$ -- масса подвижного объекта, $k_s$ -- коэффициент жесткости пружины, $f_0$ -- подъемная сила пружины в нижнем положении платформы, $F,~D$ -- нормированные проекции векторов положения шатунов в подвижной и неподвижной системе координат, $E = F - D$ -- слагаемое, характеризующее горизонтальные смещения центра масс. 

Формулы для вращающих моментов относительно рассчитанных значений $k_i$:
\begin{equation}
\label{eq:TorquesFromCrankAnglesDetailed}
\begin{aligned}
& Q_i = k_iG_i,~i=1,\dots,3 \\
& G_1 = R \left\{ L\sin{q_1}-R\sin{q_1}\cos{q_1}-R^2 \left[ \frac{1}{8b} {\left( \cos{q_1}-\cos{q_2} \right)}^2 \cos{q_1} + \frac{1}{2L}\sin^3{q_1} \right] \right. +\\ 
& \left.+ \frac{R^3}{8Lb} \left( \cos{q_1}-\cos{q_2} \right)^2 \left(\cos{2q_2}+\cos{q_1}\cos{q_2}\right) 
\right\} \\
& G_2 = R \left\{ L\sin{q_2}-R\sin{q_2}\cos{q_2}-R^2 \left[ \frac{2}{8b} {\left( \cos{q_2}-\cos{q_2} \right)}^2 \cos{q_2} + \frac{2}{2L}\sin^3{q_2} \right] \right. +\\ 
& \left.+ \frac{R^3}{8Lb} \left( \cos{q_2}-\cos{q_2} \right)^2 \left(\cos{2q_2}+\cos{q_2}\cos{q_2}\right) 
\right\} \\
& G_3 = R\left\{L\sin{q_3}-R\sin{q_3}\cos{q_3}-R^2\frac{1}{2L}\sin^3q_3\right\}
\end{aligned}
\end{equation}

На основе решения задачи динамики \ref{eq:KoefsLinearSystem2},\ref{eq:TorquesFromCrankAnglesDetailed} была получена ограничительная динамическая характеристика, проиллюстрированная на рисунке \ref{img:dyn_restrict_static}.

\begin{figure}[ht]
    \center{
        %\hfill
        \subcaptionbox[List-of-Figures entry]{\label{img:dyn_restrict_Q800}} {\includegraphics[width=0.45\linewidth]{dyn_restrict_Q800}}%
        \hfill       
        \subcaptionbox{\label{img:dyn_restrict_Q600}} {\includegraphics[width=0.45\linewidth]{dyn_restrict_Q600}}
        \hfill
    }
    \center{
        \subcaptionbox[List-of-Figures entry]{\label{img:dyn_restrict_Q500}} {\includegraphics[width=0.45\linewidth]{dyn_restrict_Q500}}
        \hfill
        \subcaptionbox[List-of-Figures entry]{\label{img:dyn_restrict_Q350}} {\includegraphics[width=0.45\linewidth]{dyn_restrict_Q350}}
        \hfill
    }
    \caption{Динамическая ограничительная характеристика 3DOF СП в декартовых координатах} % Этот текст попадает в названия рисунков в списке рисунков
    \label{img:dyn_restrict_static}
\end{figure}

Т.к. большинство параметров уравнений динамики на практике известны неточно, была произведена идентификация по следующей схеме. 

\vspace{-0pt}
Пусть по результатам эксперимента были получены $n$ стационарных положений платформы. 

Т.к. в процессе эксперимента вращающие моменты $Q_1,~Q_2,~Q_3$ измеряются, а $G_1,~G_2,~G_3$ рассчитываются по кинематическим уравнениям, вычисляется вектор $k$ путем решения системы \ref{eq:KoefsLinearSystem2} методом простых итераций. 

Далее определяются вспомогательные параметры $p_i$.. Для определения параметров используется метод наименьших квадратов с дополнительными весовыми множителями:

\begin{equation}
\label{eq:LeastSquareMethodFormulation}
\begin{aligned}
& J_c\left(x_c,y_c,z_c\right)=\sum_{i=1}^n\left[
c_{i,x}\left|\bar{p}_{3,i}-\left(\bar{S}_x^0\bar{k}\right)_iz_c + \bar{x_c}\right|^2 + c_{i,y}\left|\bar{p}_{2,i}-\left(\bar{S}_y^0\bar{k}\right)_iz_c  \bar{y_c}\right|^2\right], \\
& J_c\left(k_s,~f_0\right)=\sum_{i=1}^n{\left|\bar{p}_{1,i}-k_s\bar{h_i}-\frac{mg-f_0}{L}\right|}^2, J_c\left(x_c,~y_c,~z_c\right)\rightarrow\min, J_c\left(k_s,~f_0\right)\rightarrow\min.
\end{aligned}
\end{equation}

где $c_{i,x}$ и $c_{i,y}$  – весовые множители, $i = 1,\dots,n$ – номера стационарных участков.

При гармонических колебаний платфоры оценка моментов сил на валах кривошипов раскладывается на три компоненты:

\begin{equation}
\begin{aligned}
	m = m_0 + m_{\alpha} + m_{\beta},~m_0 = Gk_0,~m_{\alpha} = Gk_{\alpha},~m_{\beta} = Gk_{\beta}
\end{aligned}
\end{equation}

где $m_0$ -- статический момент, $m_{\alpha}$ -- момента от второй производной угла крена, $m_{\beta}$ -- момент от второй производной угла тангажа.

На каждом однородном участке колебаний с постоянной частотой это простые гармоники. При помощи БПФ определяется частота $F$ простых гармоник вращающих моментов. Величина $F$ измеряется в количестве периодов на длину выбранного участка сигнала $N$.
Далее определяются комплексные амплитуды этих гармоник:

\begin{equation}
\begin{aligned}
	m_0(t) = \operatorname{Re}\left(A_0 \exp{\frac{2\pi i}{N}Ft}\right),\\
	m_{\alpha}(t) = \operatorname{Re}\left(A_{\alpha} \exp{\frac{2\pi i}{N}Ft}\right),\\
	m_{\beta}(t) = \operatorname{Re}\left(A_{\beta} \exp{\frac{2\pi i}{N}Ft}\right).
\end{aligned}
\end{equation}

Каждая амплитуда $A_0, A_{\alpha}, A_{\beta}$ имеет 3 компоненты, соответствующие кривошипам. Аналогично выделяется первая гармоника измеряемой функции моментов на валах $m_{obs}$. Вектор её амплитуд обозначим $A_{obs}$. Неточно известны множители, с которыми линейные и угловые ускорения влияют на моменты на валах. Это масса кабины m и её моменты инерции $I_x$ и $I_y$. Однако влияние линейных ускорений $h$ на порядок меньше, чем угловых ускорений. Поэтому вместо $m$ можно подставить приблизительную оценку этой величины. Рассмотрим задачу настройки параметров $I_x$ и $I_y$. В измерениях $m_{obs}(t)$ возможны запаздывания, которые проявляются в изменении фазы в комплексной амплитуде $A_{obs}$. Поскольку смещение фазы неизвестно, то настройку будем проводить только по вещественным амплитудам гармоник. Пусть $i = 1,\ldots,n$ — номера участков сигнала с гармоническим изменением углов кривошипов. Определим функционал качества:

\begin{equation}
\begin{aligned}
	J\left(c_x,c_y\right) = \sum^n_{i=1}\left(\|A_{obs,i}\|^2 - \|A_{0,i} + c_xA_{\alpha,i} + c_yA_{\beta,i}\|^2\right)^2
\end{aligned}
\end{equation}

Этот функционал минимизируется по $\left(c_x, c_y\right)$ градиентным методом. При расчете предварительно задаются границы $c_{\min},~c_{\max}$ допустимых отклонений коэффициентов от 1.

В \underline{\textbf{четвертой главе}} приведены результаты исследования алгоритмов управления динамической платформой.

На рисунке \ref{img:SchemeControl} представлены варианты рассмотренные схемы системы управления СП.

\begin{figure}[ht]
    \center{
        \hfill
        \subcaptionbox[List-of-Figures entry]{\label{img:SchemeControlTaskSpace}} {\includegraphics[width=\linewidth]{SchemeControlTaskSpace}}%
        \hfill       
        \subcaptionbox{\label{img:SchemeControlSeparate}} {\includegraphics[width=\linewidth]{SchemeControl}}
        \hfill
    }
    
    Варианты схемы системы управления СП
    \caption{Схема системы управления СП} % Этот текст попадает в названия рисунков в списке рисунков
    \label{img:SchemeControl}
\end{figure}

Условно обозначим вариант системы управления, изображенный на рисунке \ref{img:SchemeControlSeparate} как <<система 1>>, а вариант, изображенный на рисунке \ref{img:SchemeControlTaskSpace} – <<система 2>>:

В системе 1 используются локальные контуры с ПД-регуляторами, в то время как в системе 2 осуществляется управление в пространстве положений платформы.

Для оценки качества переходных и установившихся процессов в системах 1 и 2 воспользуемся критерием минимума среднего квадрата ошибки:

\begin{equation}
\label{eq:KriteriaControlQuality}
\begin{aligned}
	J = \frac{\displaystyle \int_0^{T_{\text{уст}}} \left(c_1e_{\alpha}^2 + c_2e_{\beta}^2 + c_3e_{z_p}^2\right)dt}{T_{\text{уст}}} \rightarrow \min
\end{aligned}
\end{equation}

	где $c_1, c_2, c_3$ – нормирующие коэффициенты, $e_{\alpha}, e_{\beta}, e_{Zp}$ – ошибки регулирования по $\alpha, \beta, Zp$ соответственно.
	
Данный критерий выбран исходя из основного назначения СП в составе тренажерного комплекса -– движение по заданной траектории, формируемой математической моделью техники. Математическая модель АТ и СТ отвечает за корректный расчет параметров акселерационных воздействий на рабочем месте оператора, на основании которых строится траектория движения СП. 

К системе управления при этом предъявляются следующие требования:
\begin{itemize}
	\item минимальное запаздывание (расхождение между анимацией движения техники и движением платформы должно составлять не более 100-200 мс); 
	\item отсутствие перерегулирования;
	\item устойчивость при отработке задающих воздействий в диапазоне частот до 1 Гц.
\end{itemize}

В таблице \ref{tbl:KriteriaFreqTable} приведены результаты расчета критерия $J$ при отработке тестового движения $\alpha(t) = \alpha _0\sin(2\pi ft),~\beta(t) = \beta _0\sin(2\pi ft),~h=0$.
 
Для амплитуды колебаний  $\alpha _0 = \beta _0 = 10^o$, значение критерия $J$, удовлетворяющее перечисленным выше требованиям к системе, находится в пределах $J < 0,03$.

\begin{table} [htbp]%
    \centering
	\caption{Значение критерия для синусоидальных воздействий различной частоты}%
	\label{tbl:KriteriaFreqTable}
	\begin{tabular}{|c||c||c|}
        \hline     %%% верхняя линейка
    	Частота $f$, Гц	& Критерий $J_1$ & Критерий $J_2$\\
        \hline
    	0,1 &	0,0083 &	0,0096 \\
		0,2 &	0,0096 &	0,0109 \\
		0,3 &	0,0112 &	0,0126 \\
		0,4 &	0,0130 &	0,0143 \\
		0,5 &	0,0150 &	0,0159 \\
		0,6 &	0,0171 &	0,0174 \\
		0,7 &	0,0192 &	0,0188 \\
		0,8 &	0,0217 &	0,0205 \\
		0,9 &	0,0260 &	0,0251 \\
        \hline
	\end{tabular}%
\end{table}

Сравнивая данные таблицы \ref{tbl:KriteriaFreqTable} можно прийти к выводу, что отработка задающих воздействий в следящем режиме в обоих случаях происходит схожим образом, однако на больших частотах сигнала задания система 2 лучше компенсирует отклонения ошибки по скорости за счет влияния коэффициентов обратной связи. Система 2 также выгодно отличается большими возможностями настройки. Так, можно уменьшить коэффициенты по осям $\alpha--\beta$ для более плавной отработки низкочастотных сигналов или обнулить коэффициенты при $h$ для задействования большего объема рабочей зоны.

В \underline{\textbf{пятой главе}} описана структура программно-аппаратного комплекса и особенности программной реализации. Приведены диаграммы классов встраиваемого ПО контроллера движения СП.

В \underline{\textbf{шестой главе}} приведены результаты экспериментов, проведенных на 3DOF платформе подвижности с целью проверки корректности решения задач кинематики и динамики, определения динамических характеристик платформы.

\begin{figure}[h!]
    \centering
    \includegraphics[width=0.9\linewidth]{identificationSines}
    \caption{Тестирование модели СП в динамике}
    \label{img:identificationSines}
\end{figure}

\begin{figure}[h!]
    \centering
    \includegraphics[width=0.9\linewidth]{identificationStatic}
    \caption{Тестирование модели СП в статике}
    \label{img:identificationStatic}
\end{figure}

На рисунке \ref{img:identificationSines} показаны результаты исследования адекватности модели \ref{eq:TorquesFromCrankAnglesDetailed}. 

Пунктирными линиями показаны измеренные значения моментов, сплошными – расчетные. Максимальное запаздывание расчетных данных по фазе относительно экспериментальных составило порядка 40 градусов. 

Расхождения между моделью и экспериментом объясняются задержками в канале передачи данных между частотными преобразователями и контроллером движения, а также особенностями работы внутренних контуров регулирования АД.

В \underline{\textbf{заключении}} приведены основные результаты работы, которые состоят в следующем:
\begin{itemize}

\item проведен анализ существующих решений в области разработки систем подвижности тренажеров авто- и спецтехники. Формализованы типовые эффекты движения ТС. Сформулированы требования к системе подвижности.

\item получены уравнения обратной и прямой задач кинематики, рассмотрено решение ПЗК аналитическим и численным методами. Проведено исследование точности приближенного аналитического решения ПЗК. Была также получена зона безопасных положений платформы, вытекающая из геометрических размеров платформы. 

\item была сформулирована модель динамики системы подвижности с тремя степенями свободы и кривошипно-шатунным 
механизмом привода звеньев. Получена рабочая зона, вытекающая из паспортных значений вращающих моментов мотор-редукторов. Параметры полученной модели были уточнены по результатам эксперимента на реальном оборудовании, что позволило рассчитывать вращающие моменты по заданным движениям платформы с приемлемой точностью.  

\item приведены исследования алгоритмов управления СП. Рассмотрены два основных подхода к построению системы управления, произведено сравнение характеристик СП по критерию среднего квадрата ошибки.

\end{itemize}
Результаты работы могут быть обобщены на другие типы кинематических схем систем подвижности и использоваться в алгоритмах управления при необходимости учета динамических характеристик объекта, закрепленного на подвижном основании СП.

\newpage
\ifdefmacro{\microtypesetup}{\microtypesetup{protrusion=false}}{} % не рекомендуется применять пакет микротипографики к автоматически генерируемому списку литературы
\urlstyle{rm}                               % ссылки URL обычным шрифтом
\ifnumequal{\value{bibliosel}}{0}{% Встроенная реализация с загрузкой файла через движок bibtex8
    \renewcommand{\bibname}{\large \bibtitleauthor}
    \nocite{*}
    \insertbiblioauthor           % Подключаем Bib-базы
    %\insertbiblioexternal   % !!! bibtex не умеет работать с несколькими библиографиями !!!
}{% Реализация пакетом biblatex через движок biber
    % Цитирования.
    %  * Порядок перечисления определяет порядок в библиографии (только внутри подраздела, если `\insertbiblioauthorgrouped`).
    %  * Если не соблюдать порядок "как для \printbibliography", нумерация в `\insertbiblioauthor` будет кривой.
    %  * Если цитировать каждый источник отдельной командой --- найти некоторые ошибки будет проще.
    %
    %% authorvak
    \nocite{vakbib1}%
    \nocite{vakbib2}%
    \nocite{vakbib3}%
    %
    %% authorwos
    % \nocite{wosbib1}%
    %
    %% authorscopus
    % \nocite{scbib1}%
    %
    %% authorpatent
    % \nocite{patbib1}%
    %
    %% authorprogram
    % \nocite{progbib1}%
    %
    %% authorconf
    \nocite{confbib1}%
    % \nocite{confbib2}%
    %
    %% authorother
    \nocite{bib1}%
    \nocite{bib2}%

    \ifnumgreater{\value{usefootcite}}{0}{
        \begin{refcontext}[labelprefix={}]
            \ifnum \value{bibgrouped}>0
                \insertbiblioauthorgrouped    % Вывод всех работ автора, сгруппированных по источникам
            \else
                \insertbiblioauthor      % Вывод всех работ автора
            \fi
        \end{refcontext}
    }{
        \ifnum \totvalue{citeexternal}>0
            \begin{refcontext}[labelprefix=A]
                \ifnum \value{bibgrouped}>0
                    \insertbiblioauthorgrouped    % Вывод всех работ автора, сгруппированных по источникам
                \else
                    \insertbiblioauthor      % Вывод всех работ автора
                \fi
            \end{refcontext}
        \else
            \ifnum \value{bibgrouped}>0
                \insertbiblioauthorgrouped    % Вывод всех работ автора, сгруппированных по источникам
            \else
                \insertbiblioauthor      % Вывод всех работ автора
            \fi
        \fi
        %  \insertbiblioauthorimportant  % Вывод наиболее значимых работ автора (определяется в файле characteristic во второй section)
        \begin{refcontext}[labelprefix={}]
            \insertbiblioexternal            % Вывод списка литературы, на которую ссылались в тексте автореферата
        \end{refcontext}
        % Невидимый библиографический список для подсчёта количества внешних публикаций
        % Используется, чтобы убрать приставку "А" у работ автора, если в автореферате нет
        % цитирований внешних источников.
        \printbibliography[heading=nobibheading, section=0, env=countexternal, keyword=biblioexternal, resetnumbers=true]%
    }
}
\ifdefmacro{\microtypesetup}{\microtypesetup{protrusion=true}}{}
\urlstyle{tt}                               % возвращаем установки шрифта ссылок URL
